% frontmatter/czytam/detective.tex -- para la niña o el niño, no para el
% adulto: cómo se lee "con lupa" (las normas del club, que aparecen en la
% semana 3) y la legitymacja del Klub Lupy para rellenar el primer día.
% La de frontmatter/lupa/detective.tex, en polaco, sin género: todo en
% imperativo («Przeczytaj», «Tu narysuj»).
\section*{Jak czyta detektyw}

{\large
\begin{enumerate}[leftmargin=9mm, itemsep=3mm, label=\textbf{\arabic*.}]
\item \textbf{Przeczytaj cały tekst}, bez pośpiechu, żeby wiedzieć, o
  czym jest.
\item \textbf{Przeczytaj go jeszcze raz}, i teraz szukaj szczegółów:
  kto, gdzie, kiedy, ile, jakiego koloru, dlaczego.
\item Jeśli zdanie jest bardzo długie, \textbf{poszukaj jego części}:
  tego, co się dzieje, i słów, które to wyjaśniają -- \textit{że},
  \textit{bo}, \textit{kiedy}, \textit{chociaż}, \textit{jeśli},
  \textit{żeby}.
\item \textbf{Zanim odpowiesz, poszukaj dowodu}: zdania z tekstu, które
  przyznaje ci rację.
\item Jeśli zagadka trwa cały tydzień, \textbf{wróć do stron z
  poprzednich dni}: wskazówki nie uciekają, zostają tam.
\end{enumerate}
}

\vspace{8mm}
\begin{center}
\begin{tcolorbox}[enhanced, width=0.86\linewidth, colback=white,
  colframe=colorLectura, boxrule=1.4pt, arc=3mm,
  left=5mm, right=5mm, top=4mm, bottom=5mm,
  title={\large Klub Lupy \textperiodcentered\ Legitymacja detektywa},
  colbacktitle=colorLectura, coltitle=white, fonttitle=\bfseries\sffamily]
\begin{tabularx}{\linewidth}{@{}X r@{}}
\begin{minipage}[t]{0.56\linewidth}
\vspace{4mm}
\large
Imię: \hrulefill\par\vspace{7mm}
Wiek: \hrulefill\par\vspace{7mm}
Zwracam uwagę na: \hrulefill\par\vspace{7mm}
Podpis: \hrulefill
\end{minipage}
&
\begin{minipage}[t]{0.36\linewidth}
\vspace{2mm}
\centering
\framebox[40mm]{\rule{0pt}{44mm}}\par\vspace{1mm}
{\footnotesize\color{colorGris} Tu narysuj swój portret}
\end{minipage}
\end{tabularx}
\end{tcolorbox}

\vspace{6mm}
\iconoLupa[0.45]
\end{center}
\newpage
